% Part: computability
% Chapter: computability-theory
% Section: equiv-ce-defs

\documentclass[../../../include/open-logic-section]{subfiles}

\begin{document}

\olfileid{cmp}{thy}{eqc}

\olsection[Definisi Himpunan C. E.]{Definisi kang Setara kanggo
  Himpunan kang Bisa Dienumerasi sacara Komputabel}


Ing ngisor iki ana sawetara pratelan penting
kang setara bab teges bisa dienumerasi sacara
komputabel.

\begin{thm}
\ollabel{thm:ce-equiv}
Tetepna $S$ minangka himpunan wilangan asli.
Mula pratelan-pratelan iki setara:
\begin{enumerate}
\item\ollabel{case:ce} $S$ bisa dienumerasi sacara komputabel.
\item\ollabel{case:ran-pc} $S$ minangka citra fungsi komputabel \emph{parsial}.
\item\ollabel{case:ran-prim} $S$ kosong utawa minangka citra fungsi rekursif primitif.
\item\ollabel{case:ce-domain} $S$ minangka \emph{domain} fungsi komputabel parsial.
\end{enumerate}
\end{thm}

\begin{explain}
Telung klausa kapisan nyebut yen saben himpunan
ora kosong kang bisa dienumerasi sacara
komputabel bisa uga dienumerasi dening fungsi
komputabel, fungsi komputabel parsial, utawa
fungsi rekursif primitif. Klausa kapapat
nyebut yen $S$ bisa dienumerasi sacara
komputabel, mula kanggo sawenehing indeks~$e$,
\[
S = \Setabs{x}{\cfind{e}(x) \fdefined}.
\]
Kanthi tembung liya, $S$ yaiku himpunan
input kang ndadekake komputasi $\cfind{e}$
mandheg. Mula himpunan kang bisa dienumerasi
sacara komputabel kadhang diarani
\emph{semi-diputusake}: yen wilangan kalebu
himpunan, pungkasane kita entuk wangsulan
``iya'', nanging yen ora kalebu, kita ora
bakal entuk wangsulan ``ora''!{}
\end{explain}

\begin{proof}
Amarga saben fungsi rekursif primitif iku
komputabel lan saben fungsi komputabel uga
komputabel parsial, \olref{case:ran-prim}
ngakibatake \olref{case:ce}, dene
\olref{case:ce} ngakibatake~\olref{case:ran-pc}.
(Gatekna yen $S$ kosong, $S$~minangka citra
fungsi komputabel parsial kang ora ditegesi
ing endi wae.) Yen kita nuduhake yen
\olref{case:ran-pc} ngakibatake
\olref{case:ran-prim}, telung klausa
kapisan wis kabukten setara.

Mula upamane $S$ minangka citra fungsi
komputabel parsial $\cfind{e}$. Yen $S$
kosong, rampung. Yen ora, tetepna $a$
minangka unsur sembarang saka~$S$.
Miturut teorema wangun normal Kleene,
kita bisa nulis
\[
\cfind{e}(x) = U(\umin{s}{T(e, x, s)}).
\]
Mligine, $\cfind{e}(x) \fdefined$ lan
$= y$ yen lan mung yen ana $s$ kang
ndadekake $T(e, x, s)$ lan $U(s) = y$.
Tetepna $f(z)$ minangka
\[
f(z) = \begin{cases}
  U((z)_1) & \text{yen $T(e, (z)_0, (z)_1)$} \\
  a        & \text{yen ora.}
\end{cases}
\]
Mula $f$ rekursif primitif, amarga $T$ lan $U$
uga mangkono. \iftag{TMs}{Ing basa mesin Turing,
  yen $z$ ngode pasangan $\tuple{(z)_0, (z)_1}$
  kang $(z)_1$ minangka komputasi mesin~$M_e$
  kang mandheg kanggo input $(z)_0$, mula $f$
  ngasilake output komputasi iku; yen ora,
  ngasilake~$a$.}

Kita kudu nuduhake yen $S$ minangka citra~$f$,
yaiku kanggo saben wilangan asli~$y$,
$y \in S$ yen lan mung yen kalebu
citra~$f$. Ing arah maju, upamane $y \in S$.
Mula $y$ kalebu citra $\cfind{e}$, saengga
kanggo sawenehing $x$ lan~$s$, $T(e,x,s)$
lumaku lan $U(s) = y$. Nanging mula
$y = f(\tuple{x,s})$. Kosok baline,
upamane $y$ kalebu citra~$f$. Mula $y = a$,
utawa kanggo sawenehing~$z$,
$T(e,(z)_0,(z)_1)$ lan $U((z)_1) = y$.
% OLPL-137: Sumber beku nganggo x kang ora kaiket ing kasus iki; input sing dipilih pasangan kode yaiku (z)_0.
Amarga ing kasus kapindho $\cfind{e}((z)_0) \fdefined =
y$, ing loro-lorone kasus $y$ kalebu~$S$.

(Notasi $\cfind{e}(x) \fdefined = y$ ateges
``$\cfind{e}(x)$ ditegesi lan padha karo
$y$.'' Kita uga bisa nggunakake $\cfind{e}(x) =
y$, nanging panah tambahan kadhang migunani
kanggo ngelingake yen fungsi sing dirembug
iku parsial.)

Kanggo ngrampungake bukti \olref{thm:ce-equiv},
cukup nuduhake yen \olref{case:ce} lan
\olref{case:ce-domain} setara. Kaping pisan,
kita nuduhake yen \olref{case:ce}
ngakibatake~\olref{case:ce-domain}.
% OLPL-138: Bukti sumber ngliwati kasus S kosong; domain fungsi parsial kang ora nate ditegesi nyukupi kasus iki.
Kanggo himpunan kosong, domain fungsi komputabel
parsial kang ora nate ditegesi wis nyukupi.
Ing kasus liyane, upamane $S$ minangka citra
fungsi komputabel~$f$, yaiku
\[
S = \Setabs{y}{\text{kanggo sawenehing $x$, } f(x) = y}.
\]
Tetepna
\[
g(y) = \umin{x}{(f(x) = y)}.
\]
Mula $g$ fungsi komputabel parsial, lan
$g(y)$ ditegesi yen lan mung yen kanggo
sawenehing~$x$, $f(x) = y$. Kanthi tembung
liya, domain $g$~yaiku citra~$f$.
\iftag{TMs}{Ing basa mesin Turing: yen ana
mesin Turing~$F$ kang ngenumerasi unsur-unsur
saka~$S$, tetepna $G$ minangka mesin Turing
kang semi-mutusake $S$ kanthi nggoleki ing
antarane output~$F$ apa unsur tartamtu kalebu
himpunan; mesin iki mandheg yen unsur kasebut
ketemu, lan terus nggoleki tanpa wates yen
ora ketemu.}{}

Pungkasane, kanggo nuduhake yen
\olref{case:ce-domain} ngakibatake~\olref{case:ce},
upamane $S$~minangka domain fungsi komputabel
parsial~$\cfind{e}$, yaiku
\[
S = \Setabs{x}{\cfind{e}(x) \fdefined}.
\]
Yen $S$ kosong, rampung; yen ora,
tetepna $a$ minangka unsur sembarang
saka~$S$. Tetepna $f$ minangka
\[
f(z) = \begin{cases}
(z)_0 & \text{yen $T(e,(z)_0,(z)_1)$} \\
a & \text{yen ora.}
\end{cases}
\]
Mula, kaya ing ndhuwur, wilangan $x$
kalebu citra~$f$ yen lan mung yen
$\cfind{e}(x) \fdefined$, yaiku yen lan
mung yen $x \in S$. \iftag{TMs}{Ing basa mesin
Turing: yen ana mesin $M_e$ kang semi-mutusake
$S$, unsur-unsur $S$ bisa dienumerasi kanthi
nglakokake kabeh komputasi mesin Turing
kang mungkin, banjur ngasilake input
kang cocog karo komputasi kang mandheg.}{}
\end{proof}

Klausa~\olref{case:ce-domain} ing
\olref{thm:ce-equiv} menehi cara kang trep
kanggo ngenumerasi himpunan kang bisa
dienumerasi sacara komputabel: kanggo
saben~$e$, tetepna $W_e$ minangka domain
$\cfind{e}$, yaiku
\[
W_e = \Setabs{x}{\cfind{e}(x) \fdefined}.
\] 
Mula yen $A$ minangka himpunan sembarang
kang bisa dienumerasi sacara komputabel,
$A = W_e$ kanggo sawenehing~$e$.

Ing ngisor iki ana ciri liyane kanggo himpunan
kang bisa dienumerasi sacara komputabel.

\begin{thm}
\ollabel{thm:exists-char}
Himpunan $S$ bisa dienumerasi sacara komputabel
yen lan mung yen ana relasi komputabel
$R(x,y)$ kang ndadekake
\[
S = \Setabs{ x }{ \lexists[y][R(x,y)] }.
\]
\end{thm}

\begin{proof}
Ing arah maju, upamane $S$ bisa dienumerasi
sacara komputabel. Mula kanggo sawenehing
$e$, $S = W_e$. Kanggo nilai~$e$ iki,
kita bisa nulis $S$ minangka
\[
S = \Setabs{ x }{ \lexists[y][T(e, x, y)] }.
\]
Ing arah kosok baline, upamane
$S = \Setabs{ x }{
  \lexists[y][R(x, y)] }$. Tetepna $f$~minangka
\[
f(x) \simeq \umin{y}{R(x, y)}.
\]
Mula $f$ komputabel parsial, lan $S$
minangka domain~$f$.
\end{proof}


\end{document}
